% --- STRETCH MARGINS FOR THIS SECTION ---
\newgeometry{left=1.5cm, right=1.5cm, top=2cm, bottom=2cm}

% ==========================================
\section*{Complete Hanuman Chalisa}
\addcontentsline{toc}{section}{Complete Hanuman Chalisa}
% ==========================================

This section provides the complete and uninterrupted Hanuman Chalisa for recital in Devanagari, Telugu, and IAST Romanization.

\subsection*{Devanagari (\deva{देवनागरी})}
\addcontentsline{toc}{subsection}{\deva{देवनागरी} — \deva{श्रीगुरु चरन...}}

{ % Start formatting group
\large \linespread{1.5}\selectfont
\setcounter{chaupainum}{0} % Reset to 0 before starting

% --- Opening Dohas (Centered) ---
\begin{center}
    \begin{minipage}{0.8\linewidth}
        \centering
        \makebox[\linewidth][l]{\LARGE \deva{श्रीगुरु चरन सरोज रज}}\\
        \makebox[\linewidth][c]{\LARGE \deva{निज मनु मुकुरु सुधारि।}}\\
        \makebox[\linewidth][l]{\LARGE \deva{बरनउँ रघुबर बिमल जसु}}\\
        \makebox[\linewidth][c]{\LARGE \deva{जो दायकु फल चारि॥}}

    \end{minipage} \\[1em]
    
    \begin{minipage}{0.8\linewidth}
        \centering
        \makebox[\linewidth][l]{\LARGE \deva{बुद्धिहीन तनु जानिके }}\\
        \makebox[\linewidth][c]{\LARGE \deva{सुमिरौं पवन कुमार।}} \\ 
        \makebox[\linewidth][l]{\LARGE \deva{बल बुधि बिद्या देहु मोहिं }}\\
        \makebox[\linewidth][c]{\LARGE \deva{ हरहु कलेस बिकार॥}}
    \end{minipage}
\end{center}

\vspace{1em}

% --- Two Column Listing ---
\setlength{\columnsep}{40pt} % Increase space between the two main columns
\begin{multicols}{2}
\raggedcolumns 

% Define the command INSIDE the group to ensure it uses the correct font
\newcommand{\cp}[2]{%
    \refstepcounter{chaupainum}% Increment counter
    \begin{minipage}{\linewidth}
        \begin{tabular}{@{}p{3ex} p{0.88\linewidth}@{}}
            \small \textcolor{gray}{\arabic{chaupainum}.} & 
            \LARGE \deva{#1} \\
            & \LARGE \deva{#2}
        \end{tabular}
    \end{minipage}\\[0.8em]
}

\begin{flushleft}
\cp{जय हनुमान ज्ञान गुन सागर।}{जय कपीस तिहुँ लोक उजागर॥}
\cp{राम दूत अतुलित बल धामा।}{अंजनि पुत्र पवन सुत नामा॥}
\cp{महाबीर बिक्रम बजरंगी।}{कुमति निवार सुमति के संगी॥}
\cp{कंचन बरन बिराज सुबेसा।}{कानन कुंडल कुंचित केसा॥}
\cp{हाथ बज्र औ ध्वजा बिराजै।}{काँधे मूँज जनेऊ साजै॥}
\cp{संकर सुवन केसरी नंदन।}{तेज प्रताप महा जग बंदन॥}
\cp{बिद्यावान गुनी अति चातुर।}{राम काज करिबे को आतुर॥}
\cp{प्रभु चरित्र सुनिबे को रसया।}{राम लखन सीता मन बसिया॥}
\cp{सूक्ष्म रूप धरि सियहिं दिखावा।}{बिकट रूप धरि लंक जरावा॥}
\cp{भीम रूप धरि असुर संहारे।}{रामचंद्र के काज संवारे॥}
\cp{लाय सजीवन लखन जियाये।}{श्रीरघुबीर हरषि उर लाये॥}
\cp{रघुपति कीन्ही बहुत बड़ाई।}{तुम मम प्रिय भरतहि सम भाई॥}
\cp{सहस बदन तुम्हरो जस गावैं।}{अस कहि श्रीपति कंठ लगावैं॥}
\cp{सनकादिक ब्रह्मादि मुनीसा।}{नारद सारद सहित अहीसा॥}
\cp{जमु कुबेर दिगपाल जहाँ ते।}{कबि कोबिद कहि सके कहाँ ते॥}
\cp{तुम उपकार सुग्रीवहिं कीन्हा।}{राम मिलाय राज पद दीन्हा॥}
\cp{तुम्हरो मंत्र बिभीषन माना।}{लंकेस्वर भए सब जग जाना॥}
\cp{जुग सहस्र जोजन पर भानू।}{लील्यो ताहि मधुर फल जानू॥}
\cp{प्रभु मुद्रिका मेलि मुख माहीं।}{जलधि लांघि गये अचरज नाहीं॥}
\cp{दुर्गम काज जगत के जेते।}{सुगम अनुग्रह तुम्हरे तेते॥}
\cp{राम दुआरे तुम रखवारे।}{होत न आज्ञा बिनु पैसारे॥}
\cp{सब सुख लहै तुम्हारी सरना।}{तुम रच्छक काहू को डर ना॥}
\cp{आपन तेज सम्हारो आपै।}{तीनों लोक हांक तें कांपै॥}
\cp{भूत पिसाच निकट नहिं आवै।}{महाबीर जब नाम सुनावै॥}
\cp{नासै रोग हरै सब पीरा।}{जपत निरन्तर हनुमत बीरा॥}
\cp{संकट तें हनुमान छुड़ावै।}{मन क्रम बचन ध्यान जो लावै॥}
\cp{सब पर राम तपस्वी राजा।}{तिन के काज सकल तुम साजा॥}
\cp{और मनोरथ जो कोई लावै।}{सोइ अमित जीवन फल पावै॥}
\cp{चारों जुग परताप तुम्हारा।}{है परसिद्ध जगत उजियारा॥}
\cp{साधु सन्त के तुम रखवारे।}{असुर निकन्दन राम दुलारे॥}
\cp{अष्ट सिद्धि नौ निधि के दाता।}{अस बर दीन जानकी माता॥}
\cp{राम रसायन तुम्हरे पासा।}{सदा रहो रघुपति के दासा॥}
\cp{तुम्हरे भजन राम को पावै।}{जनम जनम के दुख बिसरावै॥}
\cp{अन्त काल रघुबर पुर जाई।}{जहाँ जन्म हरिभक्त कहाई॥}
\cp{और देवता चित्त न धरई।}{हनुमंत सेइ सर्ब सुख करई॥}
\cp{संकट कटै मिटै सब पीरा।}{जो सुमिरै हनुमत बलबीरा॥}
\cp{जै जै जै हनुमान गोसाईं।}{कृपा करहु गुरु देव की नाईं॥}
\cp{जो सत बार पाठ कर कोई।}{छूटहि बंदि महा सुख होई॥}
\cp{जो यह पढ़ै हनुमान चालीसा।}{होय सिद्धि साखी गौरीसा॥}
\cp{तुलसीदास सदा हरि चेरा।}{कीजै नाथ हृदय महँ डेरा॥}
\end{flushleft}
\end{multicols}

\vspace{1em}

% --- Concluding Doha ---
\begin{center}
    \begin{minipage}{0.8\linewidth}
        \centering
        \makebox[\linewidth][l]{\LARGE \deva{पवन तनय संकट हरन}}\\
        \makebox[\linewidth][c]{\LARGE \deva{मंगल मूरति रूप।}} \\ 
        \makebox[\linewidth][l]{\LARGE \deva{राम लखन सीता सहित}}\\
        \makebox[\linewidth][c]{\LARGE \deva{हृदय बसहु सुर भूप॥}}
    \end{minipage}
\end{center}

} % End formatting group

% --- RESTORE ORIGINAL MARGINS ---

\newpage

\subsection*{Telugu (\telu{తెలుగు})}
\addcontentsline{toc}{subsection}{\telu{తెలుగు}}

{ % Start formatting group
\large \linespread{1.5}\selectfont
\setcounter{chaupainum}{0} % Reset to 0 before starting

% --- Opening Dohas (Centered) ---
\begin{center}
    \begin{minipage}{0.8\linewidth}
        \centering
        \makebox[\linewidth][l]{\LARGE \telu{శ్రీగురు చరన సరోజ రజ}}\
        \makebox[\linewidth][c]{\LARGE \telu{నిజ మను ముకురు సుధారి।}} \\ 
        \makebox[\linewidth][l]{\LARGE \telu{బరనఉఁ రఘుబర బిమల జసు}}\
        \makebox[\linewidth][c]{\LARGE \telu{జో దాయకు ఫల చారి॥}}
    \end{minipage} \\[1em]
    
    \begin{minipage}{0.8\linewidth}
        \centering
        \makebox[\linewidth][l]{\LARGE \telu{బుద్ధిహీన తను జానికే}}\ 
        \makebox[\linewidth][c]{\LARGE \telu{సుమిరౌ పవన కుమార।}} \\ 
        \makebox[\linewidth][l]{\LARGE \telu{బల బుధి బిద్యా దేహు మోహిఁ}}\ 
        \makebox[\linewidth][c]{\LARGE \telu{హరహు కలేస బికార॥}}
    \end{minipage}
\end{center}

\vspace{1em}

% --- Two Column Listing ---
\setlength{\columnsep}{40pt} % Increase space between the two main columns
\begin{multicols}{2}
\raggedcolumns 

% Define the command INSIDE the group to ensure it uses the correct font
\newcommand{\cptelu}[2]{%
    \refstepcounter{chaupainum}% Increment counter
    \begin{minipage}{\linewidth}
        \begin{tabular}{@{}p{3ex} p{0.88\linewidth}@{}}
            \small \textcolor{gray}{\arabic{chaupainum}.} & 
            \large \telu{#1} \\
            & \large \telu{#2}
        \end{tabular}
    \end{minipage}\\[0.8em]
}

\begin{flushleft}
\cptelu{జయ హనుమాన జ్ఞాన గున సాగర।}{జయ కపీస తిహుఁ లోక ఉజాగర॥}
\cptelu{రామ దూత అతులిత బల ధామా।}{అంజని పుత్ర పవన సుత నామా॥}
\cptelu{మహాబీర బిక్రమ బజరంగీ।}{కుమతి నివార సుమతి కే సంగీ॥}
\cptelu{కంచన బరన బిరాజ సుబేసా।}{కానన కుండల కుంచిత కేసా॥}
\cptelu{హాథ బజ్ర ఔ ధ్వజా బిరాజై।}{కాఁధే మూఁజ జనేఊ సాజై॥}
\cptelu{సంకర సువన కేసరీ నందన।}{తేజ ప్రతాప మహా జగ బందన॥}
\cptelu{బిద్యావాన గునీ అతి చాతుర।}{రామ కాజ కరిబే కో ఆతుర॥}
\cptelu{ప్రభు చరిత్ర సునిబే కో రసియా।}{రామ లఖన సీతా మన బసియా॥}
\cptelu{సూక్ష్మ రూప ధరి సియహిఁ దిఖావా।}{బికట రూప ధరి లంక జరావా॥}
\cptelu{భీమ రూప ధరి అసుర సంహారే।}{రామచంద్ర కే కాజ సంవారే॥}
\cptelu{లాయ సజీవన లఖన జియాయే।}{శ్రీరఘుబీర హరషి ఉర లాయే॥}
\cptelu{రఘుపతి కీన్హీ బహుత బడాయీ।}{తుమ మమ ప్రియ భరతహి సమ భాయీ॥}
\cptelu{సహస బదన తుమ్హరో జస గావైఁ।}{అస కహి శ్రీపతి కంఠ లగావైఁ॥}
\cptelu{సనకాదిక బ్రహ్మాది మునీసా।}{నారద సారద సహిత అహీసా॥}
\cptelu{జము కుబేర దిగపాల జహాఁ తే।}{కబి కోబిద కహి సకే కహాఁ తే॥}
\cptelu{తుమ ఉపకార సుగ్రీవహిఁ కీన్హా।}{రామ మిలాయ రాజ పద దీన్హా॥}
\cptelu{తుమ్హరో మంత్ర బిభీషన మానా।}{లంకేశ్వర భఏ సబ జగ జానా॥}
\cptelu{జుగ సహస్ర జోజన పర భానూ।}{లీల్యో తాహి మధుర ఫల జానూ॥}
\cptelu{ప్రభు ముద్రికా మేలి ముఖ మాహీఁ।}{జలధి లాంఘి గయే అచరజ నాహీఁ॥}
\cptelu{దుర్గమ కాజ జగత కే జేతే।}{సుగమ అనుగ్రహ తుమ్హరే తేతే॥}
\cptelu{రామ దుఆరే తుమ రఖవారే।}{హోత న ఆజ్ఞా బిను పైసారే॥}
\cptelu{సబ సుఖ లహై తుమ్హారీ సరనా।}{తుమ రచ్ఛక కాహూ కో డర నా॥}
\cptelu{ఆపన తేజ సమ్హారో ఆపై।}{తీనోం లోక హాంక తేం కాంపై॥}
\cptelu{భూత పిసాచ నికట నహిఁ ఆవై।}{మహాబీర జబ నామ సునావై॥}
\cptelu{నాసై రోగ హరై సబ పీరా।}{జపత నిరంతర హనుమత బీరా॥}
\cptelu{సంకట తేం హనుమాన ఛుడావై।}{మన క్రమ బచన ధ్యాన జో లావై॥}
\cptelu{సబ పర రామ తపస్వీ రాజా।}{తిన కే కాజ సకల తుమ సాజా॥}
\cptelu{ఔర మనోరథ జో కోఈ లావై।}{సోఇ అమిత జీవన ఫల పావై॥}
\cptelu{చారోం జుగ పరతాప తుమ్హారా।}{హై పరసిద్ధ జగత ఉజియారా॥}
\cptelu{సాధు సన్త కే తుమ రఖవారే।}{అసుర నికన్దన రామ దులారే॥}
\cptelu{అష్ట సిద్ధి నౌ నిధి కే దాతా।}{అస బర దీన జానకీ మాతా॥}
\cptelu{రామ రసాయన తుమ్హరే పాసా।}{సదా రహో రఘుపతి కే దాసా॥}
\cptelu{తుమ్హరే భజన రామ కో పావై।}{జనమ జనమ కే దుఖ బిసరావై॥}
\cptelu{అన్త కాల రఘుబర పుర జాయీ।}{జహాఁ జన్మ హరిభక్త కహాయీ॥}
\cptelu{ఔర దేవతా చిత్త న ధరఈ।}{హనుమంత సేఇ సర్బ సుఖ కరఈ॥}
\cptelu{సంకట కటై మిటై సబ పీరా।}{జో సుమిరై హనుమత బలబీరా॥}
\cptelu{జై జై జై హనుమాన గోసాయీఁ।}{కృపా కరహు గురు దేవ కీ నాయీఁ॥}
\cptelu{జో సత బార పాఠ కర కోఈ।}{ఛూటహి బంది మహా సుఖ హోఈ॥}
\cptelu{జో యహ పఢై హనుమాన చాలీసా।}{హోయ సిద్ధి సాఖీ గౌరీసా॥}
\cptelu{తులసీదాస సదా హరి చేరా।}{కీజై నాథ హృదయ మహఁ డేరా॥}
\end{flushleft}
\end{multicols}

\vspace{1em}

% --- Concluding Doha ---
\begin{center}
    \begin{minipage}{0.8\linewidth}
        \centering
        \makebox[\linewidth][l]{\LARGE \telu{పవన తనయ సంకట హరన}}
        \makebox[\linewidth][c]{\LARGE \telu{మంగల మూరతి రూప।}} \\ 
        \makebox[\linewidth][l]{\LARGE \telu{రామ లఖన సీతా సహిత}}
        \makebox[\linewidth][c]{\LARGE \telu{హృదయ బసహు సుర భూప॥}}
    \end{minipage}
\end{center}

} % End formatting group
\newpage

\subsection*{IAST Romanization}
\addcontentsline{toc}{subsection}{IAST Romanization}

{ % Start formatting group
\large \linespread{1.5}\selectfont
\setcounter{chaupainum}{0} % Reset to 0 before starting

% --- Opening Dohas (Centered) ---
\begin{center}
    \begin{minipage}{0.8\linewidth}
        \centering
        \large \textit{śrīguru carana saroja raja, nija manu mukuru sudhāri |} \\ 
        \large \textit{baranaũ raghubara bimala jasu, jo dāyaku phala cāri ||}
    \end{minipage} \\[1em]
    
    \begin{minipage}{0.8\linewidth}
        \centering
        \large \textit{buddhihīna tanu jānike, sumiraũ pavana kumāra |} \\ 
        \large \textit{bala budhi bidyā dehu mohĩ, harahu kalesa bikāra ||}
    \end{minipage}
\end{center}

\vspace{1em}

% --- Two Column Listing ---
\setlength{\columnsep}{40pt} % Increase space between the two main columns
\begin{multicols}{2}
\raggedcolumns 

% Define the command INSIDE the group to ensure it uses the correct font
\newcommand{\cpiast}[2]{%
    \refstepcounter{chaupainum}% Increment counter
    \begin{minipage}{\linewidth}
        \begin{tabular}{@{}p{1ex} p{0.88\linewidth}@{}}
            \small \textcolor{gray}{\arabic{chaupainum}.} & 
            \large \textit{#1} \\
            & \large \textit{#2}
        \end{tabular}
    \end{minipage}\\[0.8em]
}

\begin{flushleft}
\cpiast{jaya hanumāna jñāna guna sāgara |}{jaya kapīsa tihũ loka ujāgara ||}
\cpiast{rāma dūta atulita bala dhāmā |}{añjani-putra pavana suta nāmā ||}
\cpiast{mahābīra bikrama bajaraṅgī |}{kumati nivāra sumati ke saṅgī ||}
\cpiast{kañcana barana birāja subesā |}{kānana kuṇḍala kuñcita kesā ||}
\cpiast{hātha bajra au dhvajā birājai |}{kā̃dhe mū̃ja janeū sājai ||}
\cpiast{saṅkara suvana kesarī nandana |}{teja pratāpa mahā jaga bandana ||}
\cpiast{bidyāvāna gunī ati cātura |}{rāma kāja karibe ko ātura ||}
\cpiast{prabhu caritra sunibe ko rasiyā |}{rāma lakhana sītā mana basiyā ||}
\cpiast{sūkṣma rūpa dhari siyahĩ dikhāvā |}{bikaṭa rūpa dhari laṅka jarāvā ||}
\cpiast{bhīma rūpa dhari asura sãhāre |}{rāmacandra ke kāja sãvāre ||}
\cpiast{lāya sajīvana lakhana jiyāye |}{śrīraghubīra haraṣi ura lāye ||}
\cpiast{raghupati kīnhī bahuta baṛāī |}{tuma mama priya bharatahi sama bhāī ||}
\cpiast{sahasa badana tumharo jasa gāvaĩ |}{asa kahi śrīpati kaṇṭha lagāvaĩ ||}
\cpiast{sanakādika brahmādi munīsā |}{nārada sārada sahita ahīsā ||}
\cpiast{jamu kubera digapāla jahā̃ te |}{kabi kobida kahi sake kahā̃ te ||}
\cpiast{tuma upakāra sugrīvahĩ kīnhā |}{rāma milāya rāja pada dīnhā ||}
\cpiast{tumharo mantra bibhīṣana mānā |}{laṅkesvara bhae saba jaga jānā ||}
\cpiast{juga sahasra jojana para bhānū |}{līlyo tāhi madhura phala jānū ||}
\cpiast{prabhu mudrikā meli mukha māhī̃ |}{jaladhi lāṅghi gaye acaraja nāhīṃ ||}
\cpiast{durgama kāja jagata ke jete |}{sugama anugraha tumhare tete ||}
\cpiast{rāma duāre tuma rakhavāre |}{hota na ājñā binu paisāre ||}
\cpiast{saba sukha lahai tumhārī saranā |}{tuma racchaka kāhū ko ḍara nā ||}
\cpiast{āpana teja samhāro āpai |}{tīnõ loka hā̃ka tẽ kā̃pai ||}
\cpiast{bhūta pisāca nikaṭa nahĩ āvai |}{mahābīra jaba nāma sunāvai ||}
\cpiast{nāsai roga harai saba pīrā |}{japata nirantara hanumata bīrā ||}
\cpiast{saṅkaṭa tẽ hanumāna chuṛāvai |}{mana krama bacana dhyāna jo lāvai ||}
\cpiast{saba para rāma tapasvī rājā |}{tina ke kāja sakala tuma sājā ||}
\cpiast{aura manoratha jo koi lāvai |}{soi amita jīvana phala pāvai ||}
\cpiast{cārõ juga paratāpa tumhārā |}{hai parasiddha jagata ujiyārā ||}
\cpiast{sādhu santa ke tuma rakhavāre |}{asura nikandana rāma dulāre ||}
\cpiast{aṣṭa siddhi nau nidhi ke dātā |}{asa bara dīna jānakī mātā ||}
\cpiast{rāma rasāyana tumhare pāsā |}{sadā raho raghupati ke dāsā ||}
\cpiast{tumhare bhajana rāma ko pāvai |}{janama janama ke dukha bisarāvai ||}
\cpiast{anta kāla raghubara pura jāī |}{jahā̃ janma hari-bhakta kahāī ||}
\cpiast{aura devatā citta na dharaī |}{hanumaṃta sei sarba sukha karaī ||}
\cpiast{saṅkaṭa kaṭai miṭai saba pīrā |}{jo sumirai hanumata balabīrā ||}
\cpiast{jai jai jai hanumāna gosāī̃ |}{kṛpā karahu guru deva kī nāī̃ ||}
\cpiast{jo sata bāra pāṭha kara koī |}{chūṭahi bandi mahā sukha hoī ||}
\cpiast{jo yaha paṛhai hanumāna calīsā |}{hoya siddhi sākhī gaurīsā ||}
\cpiast{tulasīdāsa sadā hari cerā |}{kījai nātha hṛdaya mahã ḍerā ||}
\end{flushleft}
\end{multicols}

\vspace{1em}

% --- Concluding Doha ---
\begin{center}
    \begin{minipage}{0.8\linewidth}
        \centering
        \large \textit{pavana tanaya saṅkaṭa harana, maṅgala mūrati rūpa |} \\ 
        \large \textit{rāma lakhana sītā sahita, hṛdaya basahu sura bhūpa ||}
    \end{minipage}
\end{center}

} % End formatting group
\restoregeometry
\newpage
